\documentclass[12pt,a4paper]{exam}
\usepackage[utf8]{inputenc}
\usepackage{amsmath,amssymb,amsfonts}
\usepackage{geometry}
\usepackage{enumitem}
\usepackage{graphicx}
\usepackage{hyperref}
\usepackage{xcolor}
\geometry{a4paper, top=2cm, bottom=2cm, left=2.5cm, right=2.5cm}
\hypersetup{colorlinks=true, linkcolor=blue, urlcolor=blue}
\renewcommand{\thequestion}{\textbf{\arabic{question}}}
\renewcommand{\questionlabel}{\thequestion.}
\pointname{ marks}
\pointsdroppedatright
\bracketedpoints
\setlength{\parindent}{0pt}
\setlength{\emergencystretch}{3em}
\allowdisplaybreaks
\newsavebox{\schmathbox}
\newcommand{\fitmath}[1]{%
  \sbox{\schmathbox}{$\displaystyle #1$}%
  \ifdim\wd\schmathbox>\linewidth
    \resizebox{\linewidth}{!}{\usebox{\schmathbox}}%
  \else
    \usebox{\schmathbox}%
  \fi}

\begin{document}

\thispagestyle{empty}
\begin{center}
  \textbf{\LARGE Diag} \\[0.3cm]
  \textbf{Time Allowed: 3 Hours} \hfill \textbf{Maximum Marks: 80} \\[0.3cm]
  \rule{\textwidth}{0.5pt}
\end{center}

\vspace{0.3cm}

\noindent\textbf{General Instructions:}
\begin{enumerate}
  \item {\normalfont\normalcolor This question paper contains 40 questions. All questions are compulsory.}
  \item {\normalfont\normalcolor Questions are grouped into sections A through E.}
  \item {\normalfont\normalcolor Use of calculators is NOT permitted.}
\end{enumerate}
\bigskip
\noindent\textbf{\large Section A: Multiple Choice \& Assertion-Reason (1 mark each)}

\begin{questions}
\question[1] If the circumference of a circle is $176\text{ cm}$, then its radius is:
\begin{enumerate}[label=(\Alph*)]
  \item (A) 14 cm
  \item (B) 21 cm
  \item (C) 28 cm
  \item (D) 35 cm
\end{enumerate}
\vspace{0.15cm}

\question[1] The area of a circle whose diameter is $14\text{ cm}$ is:
\begin{enumerate}[label=(\Alph*)]
  \item (A) 77 cm\textasciicircum{}2
  \item (B) 154 cm\textasciicircum{}2
  \item (C) 308 cm\textasciicircum{}2
  \item (D) 44 cm\textasciicircum{}2
\end{enumerate}
\vspace{0.15cm}

\question[1] The area of a quadrant of a circle with radius $r$ is:
\begin{enumerate}[label=(\Alph*)]
  \item $(A) 2\pi r$
  \item $(B) \pi r^2$
  \item $(C) \frac{1}{2} \pi r^2$
  \item $(D) \frac{1}{4} \pi r^2$
\end{enumerate}
\vspace{0.15cm}

\question[1] If the perimeter of a semi-circular protractor is $36\text{ cm}$, its diameter is:
\begin{enumerate}[label=(\Alph*)]
  \item (A) 14 cm
  \item (B) 7 cm
  \item (C) 21 cm
  \item (D) 28 cm
\end{enumerate}
\vspace{0.15cm}

\question[1] The area of a sector of a circle with radius $6\text{ cm}$ and central angle $60^\circ$ is:
\begin{enumerate}[label=(\Alph*)]
  \item $(A) 3\pi cm^2$
  \item $(B) 4\pi cm^2$
  \item $(C) 6\pi cm^2$
  \item $(D) 12\pi cm^2$
\end{enumerate}
\vspace{0.15cm}

\question[1] The mean of the first 5 natural numbers is:
\begin{enumerate}[label=(\Alph*)]
  \item (A) 2
  \item (B) 3
  \item (C) 4
  \item (D) 5
\end{enumerate}
\vspace{0.15cm}

\question[1] Find the mode of the following data: 2, 3, 4, 3, 5, 3, 6, 3.
\begin{enumerate}[label=(\Alph*)]
  \item (A) 2
  \item (B) 4
  \item (C) 3
  \item (D) 6
\end{enumerate}
\vspace{0.15cm}

\question[1] If the mean of $x, x+2, x+4, x+6, x+8$ is 24, find the value of $x$.
\begin{enumerate}[label=(\Alph*)]
  \item (A) 20
  \item (B) 22
  \item (C) 24
  \item (D) 26
\end{enumerate}
\vspace{0.15cm}

\question[1] The median of first 7 prime numbers is:
\begin{enumerate}[label=(\Alph*)]
  \item (A) 5
  \item (B) 7
  \item (C) 11
  \item (D) 13
\end{enumerate}
\vspace{0.15cm}

\question[1] For a given distribution, if Mean = 20 and Mode = 20, find the Median using the empirical relationship.
\begin{enumerate}[label=(\Alph*)]
  \item (A) 15
  \item (B) 18
  \item (C) 20
  \item (D) 25
\end{enumerate}
\vspace{0.15cm}

\question[1] Assertion (A): The area of a circle with radius $7$ cm is $154$ cm$^2$. Reason (R): The area of a circle is given by the formula $\pi r^2$.
\begin{enumerate}[label=(\Alph*)]
  \item (A) Both Assertion (A) and Reason (R) are true and Reason (R) is the correct explanation of Assertion (A)
  \item (B) Both Assertion (A) and Reason (R) are true but Reason (R) is NOT the correct explanation of Assertion (A)
  \item (C) Assertion (A) is true but Reason (R) is false
  \item (D) Assertion (A) is false but Reason (R) is true
\end{enumerate}
\vspace{0.15cm}

\question[1] Assertion (A): The circumference of a circle is $44$ cm, then its radius is $7$ cm. Reason (R): Circumference of a circle is $2\pi r$.
\begin{enumerate}[label=(\Alph*)]
  \item (A) Both Assertion (A) and Reason (R) are true and Reason (R) is the correct explanation of Assertion (A)
  \item (B) Both Assertion (A) and Reason (R) are true but Reason (R) is NOT the correct explanation of Assertion (A)
  \item (C) Assertion (A) is true but Reason (R) is false
  \item (D) Assertion (A) is false but Reason (R) is true
\end{enumerate}
\vspace{0.15cm}

\question[1] Assertion (A): The area of a semicircle of radius $r$ is $\frac{1}{2}\pi r^2$. Reason (R): A semicircle is half of a circle.
\begin{enumerate}[label=(\Alph*)]
  \item (A) Both Assertion (A) and Reason (R) are true and Reason (R) is the correct explanation of Assertion (A)
  \item (B) Both Assertion (A) and Reason (R) are true but Reason (R) is NOT the correct explanation of Assertion (A)
  \item (C) Assertion (A) is true but Reason (R) is false
  \item (D) Assertion (A) is false but Reason (R) is true
\end{enumerate}
\vspace{0.15cm}

\question[1] Assertion (A): If the perimeter of a circle is equal to that of a square, then the area of the circle is greater than the area of the square. Reason (R): The area of a circle with a given perimeter is always the maximum possible area for any plane figure.
\begin{enumerate}[label=(\Alph*)]
  \item (A) Both Assertion (A) and Reason (R) are true and Reason (R) is the correct explanation of Assertion (A)
  \item (B) Both Assertion (A) and Reason (R) are true but Reason (R) is NOT the correct explanation of Assertion (A)
  \item (C) Assertion (A) is true but Reason (R) is false
  \item (D) Assertion (A) is false but Reason (R) is true
\end{enumerate}
\vspace{0.15cm}

\question[1] Assertion (A): The length of an arc of a sector of a circle with radius $r$ and angle $\theta$ is $\frac{\theta}{360} \times 2\pi r$. Reason (R): The circumference of a circle is $2\pi r$.
\begin{enumerate}[label=(\Alph*)]
  \item (A) Both Assertion (A) and Reason (R) are true and Reason (R) is the correct explanation of Assertion (A)
  \item (B) Both Assertion (A) and Reason (R) are true but Reason (R) is NOT the correct explanation of Assertion (A)
  \item (C) Assertion (A) is true but Reason (R) is false
  \item (D) Assertion (A) is false but Reason (R) is true
\end{enumerate}
\vspace{0.15cm}

\question[1] Assertion (A): The mean of the first five natural numbers is 3. Reason (R): The mean of a set of observations is the sum of observations divided by the total number of observations.
\begin{enumerate}[label=(\Alph*)]
  \item (A) Both Assertion (A) and Reason (R) are true and Reason (R) is the correct explanation of Assertion (A)
  \item (B) Both Assertion (A) and Reason (R) are true but Reason (R) is NOT the correct explanation of Assertion (A)
  \item (C) Assertion (A) is true but Reason (R) is false
  \item (D) Assertion (A) is false but Reason (R) is true
\end{enumerate}
\vspace{0.15cm}

\question[1] Assertion (A): The empirical relationship between mean, median, and mode is $3 \text{Median} = \text{Mode} + 2 \text{Mean}$. Reason (R): For a symmetric distribution, mean, median, and mode are equal.
\begin{enumerate}[label=(\Alph*)]
  \item (A) Both Assertion (A) and Reason (R) are true and Reason (R) is the correct explanation of Assertion (A)
  \item (B) Both Assertion (A) and Reason (R) are true but Reason (R) is NOT the correct explanation of Assertion (A)
  \item (C) Assertion (A) is true but Reason (R) is false
  \item (D) Assertion (A) is false but Reason (R) is true
\end{enumerate}
\vspace{0.15cm}

\question[1] Assertion (A): The mode of the data 2, 3, 4, 3, 2, 3 is 3. Reason (R): Mode is the observation that occurs most frequently.
\begin{enumerate}[label=(\Alph*)]
  \item (A) Both Assertion (A) and Reason (R) are true and Reason (R) is the correct explanation of Assertion (A)
  \item (B) Both Assertion (A) and Reason (R) are true but Reason (R) is NOT the correct explanation of Assertion (A)
  \item (C) Assertion (A) is true but Reason (R) is false
  \item (D) Assertion (A) is false but Reason (R) is true
\end{enumerate}
\vspace{0.15cm}

\question[1] Assertion (A): If the mean of a distribution is 20 and the sum of observations is 100, then the number of observations is 5. Reason (R): The sum of observations = Mean $\times$ Number of observations.
\begin{enumerate}[label=(\Alph*)]
  \item (A) Both Assertion (A) and Reason (R) are true and Reason (R) is the correct explanation of Assertion (A)
  \item (B) Both Assertion (A) and Reason (R) are true but Reason (R) is NOT the correct explanation of Assertion (A)
  \item (C) Assertion (A) is true but Reason (R) is false
  \item (D) Assertion (A) is false but Reason (R) is true
\end{enumerate}
\vspace{0.15cm}

\question[1] Assertion (A): The median of the data 5, 8, 12, 16, 20 is 12. Reason (R): The median is always the middle value when the data is arranged in ascending order.
\begin{enumerate}[label=(\Alph*)]
  \item (A) Both Assertion (A) and Reason (R) are true and Reason (R) is the correct explanation of Assertion (A)
  \item (B) Both Assertion (A) and Reason (R) are true but Reason (R) is NOT the correct explanation of Assertion (A)
  \item (C) Assertion (A) is true but Reason (R) is false
  \item (D) Assertion (A) is false but Reason (R) is true
\end{enumerate}
\vspace{0.15cm}

\end{questions}
\bigskip
\noindent\textbf{\large Section B: Very Short Answer (2 marks each)}

\begin{questions}
\question[2] Find the area of a sector of a circle with radius $6 \text{ cm}$ if the angle of the sector is $60^\circ$. (Use $\pi = \frac{22}{7}$)
\vspace{0.15cm}

\question[2] The length of the minute hand of a clock is $14 \text{ cm}$. Find the area swept by the minute hand in $5$ minutes.
\vspace{0.15cm}

\question[2] Find the area of a quadrant of a circle whose circumference is $44 \text{ cm}$. (Use $\pi = \frac{22}{7}$)
\vspace{0.15cm}

\question[2] If the sum of the areas of two circles with radii $R_1$ and $R_2$ is equal to the area of a circle of radius $R$, prove that $R_1^2 + R_2^2 = R^2$.
\vspace{0.15cm}

\question[2] A chord of a circle of radius $10 \text{ cm}$ subtends a right angle at the center. Find the area of the corresponding minor segment. (Use $\pi = 3.14$)
\vspace{0.15cm}

\question[2] The mean of 5 observations is 15. If each observation is increased by 3, find the new mean.
\vspace{0.15cm}

\question[2] Find the class mark of the class interval 150-200.
\vspace{0.15cm}

\question[2] If the mode of a data set is 45 and the mean is 27, find the median using the empirical relationship.
\vspace{0.15cm}

\question[2] The frequency of a class is 12 and the total number of observations is 40. Find the relative frequency of this class.
\vspace{0.15cm}

\question[2] Given the data: 12, 15, 18, 12, 15, 12, 18. Find the range of the data.
\vspace{0.15cm}

\end{questions}
\bigskip
\noindent\textbf{\large Section C: Short Answer (3 marks each)}

\begin{questions}
\question[3] Prove that the angle between the two tangents drawn from an external point to a circle is supplementary to the angle subtended by the line segment joining the points of contact at the centre.
\vspace{0.15cm}

\question[3] Find the area of the shaded region in a circle of radius $14$ cm, where a sector subtends an angle of $60^{\circ}$ at the center. (Use $\pi = \frac{22}{7}$)
\vspace{0.15cm}

\question[3] The length of the minute hand of a clock is $14$ cm. Find the area swept by the minute hand in $5$ minutes.
\vspace{0.15cm}

\question[3] Find the area of a quadrant of a circle whose circumference is $44$ cm.
\vspace{0.15cm}

\question[3] A chord of a circle of radius $10$ cm subtends a right angle at the center. Find the area of the corresponding minor segment. (Use $\pi = 3.14$)
\vspace{0.15cm}

\question[3] The radii of two circles are $8$ cm and $6$ cm respectively. Find the radius of the circle having area equal to the sum of the areas of the two circles.
\vspace{0.15cm}

\end{questions}
\bigskip
\noindent\textbf{\large Section D: Long Answer (4-5 marks each)}

\begin{questions}
\question[5] A chord $AB$ of a circle of radius $15 \text{ cm}$ subtends an angle of $60^\circ$ at the centre. Find the areas of the corresponding minor and major segments of the circle. (Use $\pi = 3.14$ and $\sqrt{3} = 1.73$)
\vspace{0.15cm}

\question[5] An umbrella has 8 ribs which are equally spaced. Assuming the umbrella to be a flat circle of radius $45 \text{ cm}$, find the area between the two consecutive ribs of the umbrella.
\vspace{0.15cm}

\question[5] Find the area of the shaded region in a square $\text{ABCD}$ of side $14 \text{ cm}$, where four circles are drawn with centers at the four vertices of the square, each having a radius of $7 \text{ cm}$.
\vspace{0.15cm}

\question[5] A round table cover has six equal designs as shown in the figure. If the radius of the cover is $28 \text{ cm}$, find the cost of making the designs at the rate of $\text{ \text{Rs.~} } 0.35 \text{ per cm}^2$. (Use $\sqrt{3} = 1.7$)
\vspace{0.15cm}

\end{questions}
\bigskip

\end{document}
